\scanpage{24}
\uline{Lemme 2} Soit $\deleted{\unsure{A}}B$ un anneau $1/2$ local, $\deleted{\unsure{B}}C$ une algèbre \deleted{finie}
sur $B$, telle que $C$ soit $1/2$ local, et les idéaux
induits par les idéaux maximaux \inserted{\deleted{\unsure{$\mathfrak m$}} $\mathfrak n$} de $C$ sont des
idéaux maximaux \deleted{\unsure{de}} \inserted{$\mathfrak m$} de $B$. Supposons l'une
des conditions suivantes satisfaite

(i) Les corps $k(\mathfrak m)$ sont infinis

(ii) Deux idéaux maximaux \inserted{$\mathfrak n,\mathfrak n'$} distincts de $C$
induisent des idéaux maximaux \inserted{$\mathfrak m,\mathfrak m'$} distincts de $B$,
et pour tt $\mathfrak n$, l'extension résiduelle $k(\mathfrak n)/k(\mathfrak m)$
est radicielle (p. ex. $C$ fini et radiciel sur $B$
$1/2$ local).

Alors pour tt $n\geq1$, et tout élément $u\in\mathbb M_n(C)$,
$\exists\,v\in\mathbb M_n(B)$ tel que $u+\varphi(v)\in\mathbb M_n(C)^*$.

\uline{Dém}
Soient $\mathfrak m_i$ $(1\leq i\leq\unsure{n|m})$ les idéaux maximaux de
$B$ induits par ceux de $C$. Ils sont en nb fini,
donc pour tt\textsuperscript{e} famille d'éléments $v_{\unsure{i|j}}\in\mathbb M_n(k(\mathfrak m_i))$
il existe un $v\in\mathbb M_n(B)$ qui les donne. D'ailleurs,
pour un $w\in\mathbb M_n(C)$, le fait que $w$ soit
inversible ne dépend que de l'inversibilité de $\det w$,
donc que de son image dans $\mathbb M_n(C/\operatorname{rad} C)$! Ces
deux remarques nous ramènent au cas où $C$ est
un composé fini de corps, et $B$ aussi ; On
peut $\overline{\mathrm m}$ supposer $B$ un corps, $B=k$.
% archive sheet number: 23
